\section*{一、这本手册是什么}

这是一本\emph{严格按《考研数学二考试大纲》顺序}编排的高等数学知识点手册。它不替代教材与习题册，
解决的是三件事：

\begin{enumerate}
  \item \textbf{对照考纲查漏}：大纲上出现的每一条，本册都写到；大纲外的内容一律不收，避免把时间花在
        数一专属的级数、空间解析几何、三重积分与曲线曲面积分上。
  \item \textbf{按星级分配精力}：每个考点都标了五级星标，先攻$\bigstar\bigstar\bigstar\bigstar\bigstar$，再补$\bigstar\bigstar\bigstar\bigstar$，不要平均用力。
  \item \textbf{考前快速回顾}：公式、条件、易错点集中成块，配合每章末尾的易错清单，可在两小时内过完一章。
\end{enumerate}

\section*{二、五星分级标准}

星标取\textbf{考频 $\times$ 难度 $\times$ 分值}的综合，不是单纯的"重要程度"：

\begin{center}
\begin{tabular}{@{}clll@{}}
\toprule
\textbf{星标} & \textbf{名称} & \textbf{含义} & \textbf{复习要求} \\
\midrule
$\bigstar\bigstar\bigstar\bigstar\bigstar$ & 决胜考点 & 几乎年年必考，且是难题/易错点 & 闭卷默写 + 限时算对 \\
$\bigstar\bigstar\bigstar\bigstar$ & 高频重点 & 反复考，套路固定但计算量不小 & 熟练到条件反射 \\
$\bigstar\bigstar\bigstar$ & 常考点 & 隔年出现或以小题形式考 & 理解 + 会用 \\
$\bigstar\bigstar$ & 一般考点 & 低频，或只是解题的基础工具 & 会算即可 \\
$\bigstar$ & 了解性考点 & 大纲要求"了解"，考到即送分 & 记住结论 \\
\bottomrule
\end{tabular}
\end{center}

每章开头的"星级总表"给出该章全部考点的星级、大纲要求层次、难度与常见题型，可以先看表再决定精读哪几节。

\section*{三、试卷结构与分值}

数学二满分 150 分，考试时间 180 分钟，闭卷笔试。内容结构：\textbf{高等数学约 80\%}、
\textbf{线性代数约 20\%}，不考概率论与数理统计。题型结构：

\begin{center}
\begin{tabular}{@{}lll@{}}
\toprule
\textbf{题型} & \textbf{题量} & \textbf{分值} \\
\midrule
单项选择题 & 10 小题 & 每小题 5 分，共 50 分 \\
填空题 & 6 小题 & 每小题 5 分，共 30 分 \\
解答题（包括证明题） & 6 小题 & 共 70 分 \\
\bottomrule
\end{tabular}
\end{center}

由此可得一条实用结论：\textbf{选择填空共 80 分，占比过半}，而它们考的是基本概念、基本计算与常见陷阱。
本册中标注$\bigstar\bigstar\bigstar$以上、题型为"选择/填空"的考点，是性价比最高的得分点。

\section*{四、数学二不考什么（去糟粕清单）}

以下内容属于数一/数三或大纲外，本册一律不收，复习时可以直接跳过：

\begin{itemize}
  \item 向量代数与空间解析几何（平面、直线、曲面、曲线）
  \item 无穷级数（数项级数、幂级数、傅里叶级数）
  \item 三重积分、曲线积分、曲面积分、高斯公式、斯托克斯公式、场论
  \item 含参变量积分、$\Gamma$ 函数等分析类拓展
  \item 微分方程的欧拉方程、方程组、级数解法、拉普拉斯变换、差分方程
  \item 概率论与数理统计（数学二完全不考）
\end{itemize}

在线性代数部分，数学二与数一范围基本一致（行列式、矩阵、向量、线性方程组、特征值与特征向量、
二次型六章全考），但\textbf{不考}线性空间与坐标变换、线性变换、若尔当标准形等更深的定理。

\section*{五、怎么用：三轮法}

\begin{enumerate}
  \item \textbf{第一轮（打基础）}：按章节顺序读，重点看"定义与结论"和公式的\textbf{适用前提}，
        每个$\bigstar\bigstar\bigstar\bigstar$以上的考点都动手算一遍例题。
  \item \textbf{第二轮（补漏洞）}：只看每章的"星级总表"和"易错清单"，把$\bigstar\bigstar$以下的快速扫过，
        把$\bigstar\bigstar\bigstar\bigstar$以上的重新推一遍；错题回到对应考点的warning框。
  \item \textbf{第三轮（考前）}：只看$\bigstar\bigstar\bigstar\bigstar\bigstar$考点与公式，配合本册的公式速查部分，
        限时做真题，检查是否还有"知道套路但算错"的考点。
\end{enumerate}

\section*{六、使用提醒}

\begin{itemize}
  \item 本册按大纲顺序编排，因此\textbf{章号即大纲条目号}（如"一元函数微分学"对应大纲第二部分）。
  \item 手册中的"真题考法"只描述\textbf{出题角度}，不编造年份与题号；确需核实的年份，请以真题册为准。
  \item 所有公式都以考纲与统考习惯写法为准；不同教材符号不一致处，正文中已注明"统考按 \dots 理解"。
  \item 手册是复习的\textbf{索引与骨架}，不能替代刷题。看完一个考点，请立刻用对应真题或习题验证。
\end{itemize}
